\section*{Capítulo \ref{ch:b1:elementary-steps} --- Mecanismos de reacción: etapas elementales y aproximaciones}

\begin{solution}{exo:b1:elementary-steps:1}
(a) Unimolecular: puede ser elemental. (b) Bimolecular: puede ser
elemental. (c) Tres moléculas y cuatro enlaces reconstruidos a la vez: no
es elemental. (d) Trimolecular, con \ce{N2} llevándose la energía
liberada: una \omterm{def:b1:elementary-steps:elementary-step}{etapa elemental} rara, pero auténtica.
\end{solution}

\begin{solution}{exo:b1:elementary-steps:2}
$v = k[\mathrm{A}][\mathrm{B}]$; $v = k[\mathrm{A}]^2$; $v = k[\mathrm{A}]$.
Una ecuación global es un balance, no una descripción de cómo se
encuentran las moléculas; su \omterm{def:b1:rate-laws:order}{ley de velocidad} debe medirse o deducirse de
un mecanismo.
\end{solution}

\begin{solution}{exo:b1:elementary-steps:3}
El intermedio es el hueco situado a 30; los \omterm{def:b1:elementary-steps:energy-profile}{estados de transición} son los
máximos situados a 70 y 55. El \omterm{def:b1:elementary-steps:energy-profile}{estado de transición} más alto es el
primero: la primera etapa es la determinante. Es endotérmica (de 0 sube a
30).
\end{solution}

\begin{solution}{exo:b1:elementary-steps:4}
Por el mismo factor, $10^4$ (\cref{prop:b1:elementary-steps:catalysis});
la constante de equilibrio no cambia; el equilibrio se alcanza unas
$10^4$ veces antes.
\end{solution}

\begin{solution}{exo:b1:elementary-steps:5}
$t_{\max} = \ln(0.30/0.10)/(0.30 - 0.10) = \qty{5.49}{s}$. Entonces,
$[\mathrm{B}]_{\max} = a_0\frac{0.10}{0.20}(\eu^{-0.549} - \eu^{-1.648}) =
0.5\,a_0(0.577 - 0.192) = 0.192\,a_0$.
\end{solution}

\begin{solution}{exo:b1:elementary-steps:6}
$\dd[\mathrm{I}]/\dd t = k_1[\mathrm{A}] - (k_{-1} + k_2)[\mathrm{I}] = 0$
da $[\mathrm{I}] = k_1[\mathrm{A}]/(k_{-1} + k_2)$ y $v = k_2[\mathrm{I}]
= \frac{k_1k_2}{k_{-1}+k_2}[\mathrm{A}]$. Si $k_2 \gg k_{-1}$, $v =
k_1[\mathrm{A}]$: la primera etapa es la determinante. Si $k_2 \ll k_{-1}$,
$v = (k_1/k_{-1})k_2[\mathrm{A}]$: un preequilibrio seguido de una etapa
lenta.
\end{solution}

\begin{solution}{exo:b1:elementary-steps:7}
Suma: \ce{2H2O2 -> 2H2O + O2}. \omterm{def:b1:elementary-steps:catalyst}{Catalizador}: \ce{I-} (consumido y después
regenerado); intermedio: \ce{IO-}. La primera etapa, lenta, fija la
velocidad: $v = k_1[\ce{H2O2}][\ce{I-}]$.
\end{solution}

\begin{solution}{exo:b1:elementary-steps:8}
La etapa es muy cuesta arriba: por el postulado de Hammond, su \omterm{def:b1:elementary-steps:energy-profile}{estado de
transición} se parece al carbocatión. Un sustituyente que rebaja la
energía del carbocatión rebaja casi tanto la del \omterm{def:b1:elementary-steps:energy-profile}{estado de transición}, de
modo que la barrera baja y la etapa es más rápida.
\end{solution}

\begin{solution}{exo:b1:elementary-steps:9}
$k_{\mathrm B}/k_{\mathrm C} = \exp(20000/RT)$: \num{3.0e3} a \qty{300}{K},
$\exp(4.01) = 55$ a \qty{600}{K}. Calentar reduce la preferencia cinética
por B y, al hacer reversible la formación de B, permite que el sistema
alcance C, más estable.
\end{solution}

\begin{solution}{exo:b1:elementary-steps:10}
$k_1[\mathrm{A}][\mathrm{M}] = (k_{-1}[\mathrm{M}] + k_2)[\mathrm{A}^*]$,
de modo que $v = k_2[\mathrm{A}^*] = \dfrac{k_1k_2[\mathrm{A}][\mathrm{M}]}{k_{-1}[\mathrm{M}]
+ k_2}$. A presión alta ($k_{-1}[\mathrm{M}] \gg k_2$), $v =
(k_1k_2/k_{-1})[\mathrm{A}]$: orden 1. A presión baja, $v =
k_1[\mathrm{A}][\mathrm{M}]$: orden 2; la activación por choque se
convierte en la etapa lenta.
\end{solution}

\begin{solution}{exo:b1:elementary-steps:11}
La demostración es la del teorema. Para $k_1 = k_2 = k$:
$\frac{\dd}{\dd t}([\mathrm{B}]\eu^{kt}) = ka_0$, de modo que $[\mathrm{B}]\eu^{kt} =
ka_0t$ y $[\mathrm{B}] = a_0kt\,\eu^{-kt}$, máxima en $t = 1/k$.
\end{solution}

\begin{solution}{exo:b1:elementary-steps:12}
$\mathrm{A} + \mathrm{B} \rightleftharpoons \mathrm{I} + \mathrm{C}$
(rápida, constante $K_1$) y, después, $\mathrm{I} + \mathrm{B} \to \mathrm{P}$
(lenta, $k_2$). La suma es la reacción global; $[\mathrm{I}] =
K_1[\mathrm{A}][\mathrm{B}]/[\mathrm{C}]$ y $v = k_2[\mathrm{I}][\mathrm{B}]
= k_2K_1[\mathrm{A}][\mathrm{B}]^2/[\mathrm{C}]$: el subproducto, liberado
en el preequilibrio, frena la reacción.
\end{solution}

\begin{solution}{pb:b1:elementary-steps:1}
\textbf{1.} Duplicar $[\ce{NO}]_0$ multiplica $v_0$ por 4: orden 2;
duplicar $[\ce{O2}]_0$ la multiplica por 2: orden 1.
\textbf{2.} $v = k[\ce{NO}]^2[\ce{O2}]$; $k = 1.0 \times 10^{-5}/((1.0
\times 10^{-3})^2 \times 1.0 \times 10^{-3}) =
\qty{1.0e4}{L^2.mol^{-2}.s^{-1}}$.
\textbf{3.} Un choque simultáneo de tres moléculas es raro, y una sola
etapa no podría hacerse más lenta al calentar (pregunta 5).
\textbf{4.} $E_a = R\ln(k_{400}/k_{300})/(1/300 - 1/400) = 8.314 \ln(0.1)
/(8.33 \times 10^{-4}) = \qty{-23}{kJ/mol}$.
\textbf{5.} Una \omterm{def:b1:elementary-steps:elementary-step}{etapa elemental} transcurre mediante choques que superan
una barrera; la fracción de ellos capaces de hacerlo, $\eu^{-E_a/RT}$ con
$E_a \ge
0$, solo puede crecer con la temperatura.
\textbf{6.} $-\dd[\ce{NO}]/\dd t = 2v$, $-\dd[\ce{O2}]/\dd t = v$.
\textbf{7.} \ce{2NO -> N2O2} más \ce{N2O2 + O2 -> 2NO2} da
\ce{2NO + O2 -> 2NO2}; el intermedio es \ce{N2O2}.
\textbf{8.} Ambas bimoleculares (y la inversa de la primera,
unimolecular).
\textbf{9.} $v = k_2[\ce{N2O2}][\ce{O2}]$.
\textbf{10.} $[\ce{N2O2}] = K_1[\ce{NO}]^2$ con $K_1 = k_1/k_{-1}$.
\textbf{11.} $v = k_2K_1[\ce{NO}]^2[\ce{O2}]$: la ley observada.
\textbf{12.} $k = k_1k_2/k_{-1}$.
\textbf{13.} $\dd[\ce{N2O2}]/\dd t = k_1[\ce{NO}]^2 - k_{-1}[\ce{N2O2}] -
k_2[\ce{N2O2}][\ce{O2}]$.
\textbf{14.} Al igualarla a cero: $[\ce{N2O2}] = k_1[\ce{NO}]^2/(k_{-1} +
k_2[\ce{O2}])$.
\textbf{15.} $v = k_2[\ce{N2O2}][\ce{O2}]$ da la ley enunciada.
\textbf{16.} $k_{-1} \gg k_2[\ce{O2}]$: el dímero se rompe mucho más a
menudo de lo que se encuentra con el dioxígeno, de modo que la primera
etapa permanece en equilibrio.
\textbf{17.} Si $k_2[\ce{O2}] \gg k_{-1}$, $v = k_1[\ce{NO}]^2$: orden 2
respecto de \ce{NO} y 0 respecto de \ce{O2}.
\textbf{18.} El orden 1 medido respecto de \ce{O2}: el primer límite.
\textbf{19.} $\ln k = \ln k_1 + \ln k_2 - \ln k_{-1}$.
\textbf{20.} Al derivar respecto de $T$ y usar $\dd\ln k_i/\dd T
= E_i/RT^2$: $E_a = E_1 + E_2 - E_{-1}$.
\textbf{21.} Calentar acelera un poco la etapa lenta ($E_2$ pequeña), pero
desplaza mucho el preequilibrio hacia \ce{NO} ($E_{-1}$ grande): hay
menos dímeros, y la velocidad neta disminuye.
\textbf{22.} El dímero está $E_1 - E_{-1} = \qty{-38}{kJ/mol}$ por debajo
de los reactivos, y el segundo \omterm{def:b1:elementary-steps:energy-profile}{estado de transición}, $E_2 = \qty{15}{kJ/mol}$
por encima del dímero: \qty{23}{kJ/mol} por debajo de los reactivos. El
punto más alto del camino es el primer \omterm{def:b1:elementary-steps:energy-profile}{estado de transición}, solo
\qty{2}{kJ/mol} más arriba.
\textbf{23.} $\exp\big(\frac{23000}{8.314}(\frac{1}{400} - \frac{1}{300})\big)
= 0.10$: diez veces más lenta a \qty{400}{K}.
\textbf{24.} $E_a = 2 - 40 + 15 = \textbf{\qty{-23}{kJ/mol}}$, el valor
medido en la parte I: el mecanismo explica a la vez la \omterm{def:b1:rate-laws:order}{ley de velocidad} y
su extraña dependencia de la temperatura.
\end{solution}
